\subsection{Removing the Gram cutoff at level zero}
\label{actual-gram-cutoff-pbw}

The cutoff can be removed without estimating differentiated terminal
word densities.  The reason is that the original level-zero form is a
Jacquet pairing with a product of central Capelli operators.  Compact
invariance of the Gram polynomial and a relative
Poincare--Birkhoff--Witt filtration bound the
number of derivatives which can hit the cutoff independently of the
compact dimension.  This gives a sufficient quadratic Gram power;
it does not assert the sharper source statement for every smaller
admissible power.

Put $s=r+1$, $n=k+2r+1$, and retain the original Hermitian space
$V=U\oplus W\oplus U^\vee$, with $\dim U=r$ and $W$ of
signature $(k,1)$.  The original right compact group
$K=K_{z'}$ fixes $U,U^\vee,z'$ and acts unitarily on
$W'_+=W\cap(z')^\perp$.  Thus $K\simeq U(k)$.
It is this particular subgroup which is used below.  Its invariance
is required only of the Gram polynomial.  The input coefficient
need not be finite under this compact group or covariant for a
specified compact type.

Write $\mathcal G$ for the full-group Gram polynomial.  The proved
cofactor construction gives polynomial coordinates $\Omega_j$, each
of degree $s$, such that
\begin{equation}\label{pbw-gram-squares}
 \mathcal G=T'\sum_{j=1}^k|\Omega_j|^2,\qquad T'>0.
\end{equation}
The identity holds on the entire group, including the zero set.
Both $\mathcal G$ and all original leading Bessel minors are
right $K$-invariant, because $K$ fixes the vectors defining their
columns.  The decay seminorms $\|\cdot\|_{A,m}$ and the derivatives
\[
 L_XF(g)=\left.\frac d{dt}F(e^{-tX}g)\right|_0,
 \qquad R_XF(g)=\left.\frac d{dt}F(ge^{tX})\right|_0
\]
are those of the established Jacquet pairing theorem, named
\emph{Equality at every positive character} in the construction
of regularization, unfolding, and Jacquet adjointness.  Let $B$ be
any exponent for which that theorem gives
\begin{equation}\label{pbw-pairing-bound}
 |\beta(F,H)|\le C_B\|F\|_{B,4}\|H\|_{B,4}
 \quad(F,H\in\mathcal C_B^5).
\end{equation}
The original finite-decay reserve supplies such a $B$.

\paragraph{The actual level-zero identity.}
For an original admissible word $w=((\epsilon_1,t_1),\ldots,
(\epsilon_L,t_L))$, put
\[
 \mathcal C_w=\prod_{j=1}^L\sigma\mathcal C_V(t_j),
 \qquad \sigma=(-1)^{k-l_0},
\]
where $\mathcal C_V(t)$ is the original full Capelli operator,
with parameter $t+(n-1)/2$.  All factors commute.  Both signs of
letter use their actual numerical parameter $t_j$; in particular a
negative grid letter retains its shift $y+1$.  The original forms
obey
\begin{equation}\label{pbw-level-zero-identity}
 \operatorname{RR}_{0,w}(F;\varnothing,\varnothing)
                    =\beta(\mathcal C_wF,F).
\end{equation}
Here is the reduction, including the word corner.  At level zero
the bit-zero letter is the identity, the bit-one letter is zero,
and the empty-head slice is the identity.  Thus only the all-zero
auxiliary index contributes to the form.  In the original residue
recursion, an unused tensor coordinate is free in the old residue;
the first column of the dressed $00$ transfer entry at that index
is the bare $Q_{00}$ times the all-zero vector.  All other columns
are killed by freeness at the new coordinate.  Consequently
\[
 (R_{l::w})_{00}=(R_w)_{00}Q_{00}(t_l).
\]
The same equation holds for a negative letter since transposition
does not change the $00$ entry.  Induction, with the empty residue
equal to the identity, gives the reversed product of $Q_{00}$'s.
At level zero $Q_{00}(t)=\sigma\mathcal C_V(t)$, so central
commutativity removes the reversal and proves
\eqref{pbw-level-zero-identity}.  This is a finite algebraic
reduction of the prescribed residue form; no terminal integral
limit is used.

\begin{lemma}[Derivatives which can hit a compact-invariant multiplier]
\label{pbw-relative-order}
Let $q=2(n-k)=4s-2$.  For every word of length $L$, the original
operator $\mathcal C_w$ is a finite sum
\begin{equation}\label{pbw-reduced-word}
 \mathcal C_w=\sum_\alpha c_\alpha
 R_{X_{\alpha,1}}\cdots R_{X_{\alpha,h_\alpha}}
 R_{Z_{\alpha,1}}\cdots R_{Z_{\alpha,v_\alpha}},
 \qquad h_\alpha\le qL,\quad h_\alpha+v_\alpha\le nL,
\end{equation}
where $Z_{\alpha,j}\in\mathfrak k_{\mathbb C}$ and all Lie
directions and scalar coefficients are constant.  Consequently,
for every smooth right $K$-invariant scalar function $\eta$,
\[
 \mathcal C_w(\eta f)=\sum_\alpha c_\alpha
 R_{X_{\alpha,1}}\cdots R_{X_{\alpha,h_\alpha}}
 \bigl(\eta\,R_{Z_{\alpha,1}}\cdots
                      R_{Z_{\alpha,v_\alpha}}f\bigr).
\]
In particular at most $qL$ derivatives of this expression can hit
$\eta$, although as many as $nL$ may hit $f$.
\end{lemma}

\begin{proof}
Choose a basis adapted to $W'_+$ and its complement.  The
complexified Lie algebra of $K$ is the $k\times k$ block
$\mathfrak k_{\mathbb C}=\mathfrak{gl}_k$ inside
$\mathfrak{gl}_n$.  Let $\mathfrak p_0$ be the span of matrix
units with at least one index outside this block.  Then
$[\mathfrak k_{\mathbb C},\mathfrak p_0]\subset\mathfrak p_0$.
Give a word degree equal to its number of $\mathfrak p_0$ factors;
compact factors have degree zero.

In the original source convention, the generators are
$E_{ab}=R_{h_1(E_{ab}^{\rm mat})}
       +iR_{h_2(E_{ab}^{\rm mat})}$, where
$h_1(Z)=(Z-H^{-1}Z^*H)/2$ and
$h_2(Z)=(Z+H^{-1}Z^*H)/(2i)$.  Thus they are precisely the
complexified right derivatives; $h_1(Z)+ih_2(Z)=Z$.
The full Capelli determinant is independent of the chosen basis,
and in this basis it is its usual column determinant with scalar
diagonal shifts.  In any permutation term, at most $n-k$ factors
have their column outside the compact block, and at most another
$n-k$ have their row outside it.  Thus every term contains at most
$q$ noncompact factors.  Scalar choices from the diagonal shifts
can only reduce that count.  Expanding a product of $L$ such
determinants gives degree at most $qL$.

Move all compact Lie algebra factors to the right using
$XY=YX+[X,Y]$.  A commutator of a compact factor and an outside
factor is again outside, so the degree never increases.  PBW,
or induction on the number of misplaced compact factors and word
length, terminates this process.  Every term is an outside word
of length at most $qL$ followed by a word in
$U(\mathfrak k_{\mathbb C})$.  The original total length is
at most $nL$ and every commutator reduces total length, proving
both bounds in \eqref{pbw-reduced-word}.  Since $R_Z\eta=0$
for every $Z\in\mathfrak k_{\mathbb C}$, the compact suffix
commutes with multiplication by $\eta$ and acts only on $f$.

The argument is applied to the whole product.  It does not assume
that an individual quantum minor is central, that intermediate
outputs preserve a fixed type, or that the coefficient itself
has finite compact type.  The right derivatives have the
convention $g\exp(tX)$, and the compact directions are those of
the actual $K_{z'}$.  Thus their annihilation of $\eta$ is exact.
If one instead starts with the left convention
$L_XF(g)=dF(e^{-tX}g)/dt|_0$, the corresponding principal
anti-automorphism reverses words and changes the indicated signs.
Both operations preserve the outside-factor filtration, so the
same order bound holds; no equality of naively ordered left and
right generators is asserted.
\end{proof}

\begin{lemma}[Uniform cutoff error]
\label{pbw-cutoff-error}
Let $\chi\in C^\infty(\mathbb R)$ equal zero on $[0,1/2]$
and one on $[1,\infty)$, and let $p,m\ge0$ be integers with
$2p>m$.  For $0<\epsilon\le1$ and every smooth $f$,
\begin{equation}\label{pbw-error-estimate}
 \|\bigl(\chi(\mathcal G/\epsilon)-1\bigr)
                \mathcal G^pf\|_{B,m}
 \le C\epsilon^{p-m/2}\|f\|_{B+sm,m}.
\end{equation}
The constant is independent of $\epsilon$.
\end{lemma}

\begin{proof}
On the real vector space underlying $\mathbb C^k$, set
\[
 h_\epsilon(v)=
 (\chi(T'\|v\|^2/\epsilon)-1)(T'\|v\|^2)^p.
\]
The function $h_1$ is smooth and compactly supported, including
at $v=0$, and
$h_\epsilon(v)=\epsilon^ph_1(v/\sqrt\epsilon)$.
Its derivative of order $a$ is therefore bounded by
$C_a\epsilon^{p-a/2}$, uniformly in $v$.
Every invariant derivative of every $\Omega_j$ is a polynomial
of degree at most $s$: invariant differentiation is linear on
matrix entries and does not increase degree.  The ordered chain
rule applied to $h_\epsilon(\Omega(g))$ thus bounds any mixed
left/right derivative of total length $j\le m$ by
\[
 C_j\epsilon^{p-j/2}M(g)^{sj}.
\]
Here $M(g)\ge1$ on the group, and $\epsilon\le1$ permits
replacing every smaller chain-rule order by $j$.  The product
rule with $f$ proves \eqref{pbw-error-estimate}.  In particular,
the estimate covers the singular set $\mathcal G=0$; no local
extension across that set is being assumed.
\end{proof}

\begin{corollary}\label{pbw-cutoff-word-error}
Put $m=5+qN$ and $J=5+nN$.  If $2p>m$ and $|w|\le N$, then
\[
 \|\mathcal C_w[(\chi(\mathcal G/\epsilon)-1)
                         \mathcal G^pf]\|_{B,5}
 \le C_w\epsilon^{p-m/2}\|f\|_{B+sm,J}
                 \qquad(0<\epsilon\le1).
\]
\end{corollary}

\begin{proof}
In each term of \eqref{pbw-reduced-word}, the compact suffix
acts only on $f$.  Apply Lemma~\ref{pbw-cutoff-error} with
order $h_\alpha+5$ to that compact derivative of $f$.
At most $h_\alpha+5\le m$ derivatives hit the scalar error,
and the total number hitting $f$ is at most
$h_\alpha+v_\alpha+5\le J$.  Every complexified Lie direction
is a fixed complex linear combination of real directions, so
the original real derivative seminorms control these terms.
Summing the finite expansion proves the estimate.
\end{proof}

\begin{theorem}[Level-zero cutoff convergence]
\label{pbw-level-zero-cutoff}
Fix $N\ge0$, put $m=5+(4s-2)N$, and choose an integer $p$ with
$2p>m$.  Suppose $f$ is smooth with all finite derivative
seminorms at one decay exponent $A$, and
$A\ge B+2sp$.  Put
\[
 F=\mathcal G^pf,\qquad
 F_\epsilon=\chi(\mathcal G/\epsilon)\mathcal G^pf.
\]
For every original admissible word $|w|\le N$,
\begin{equation}\label{pbw-level-zero-limit}
 \operatorname{RR}_{0,w}(F_\epsilon;\varnothing,\varnothing)
 \longrightarrow
 \operatorname{RR}_{0,w}(F;\varnothing,\varnothing)
                    \qquad(\epsilon\downarrow0).
\end{equation}
The same conclusion holds for the original normalized energies,
with their original fixed positive normalizers.
For $N=3s$, the choice $p=6s^2$ satisfies the required inequality
in every rank, including $s=1$.
\end{theorem}

\begin{proof}
Since $sm<2sp$, the hypothesis supplies
$\|f\|_{B+sm,J}<\infty$ for $J=5+nN$.
Lemma~\ref{pbw-cutoff-error} gives $F_\epsilon\to F$ in
$\mathcal C_B^5$, and Corollary~\ref{pbw-cutoff-word-error}
gives $\mathcal C_wF_\epsilon\to\mathcal C_wF$ in the same
space.  The uncut $F$ and $\mathcal C_wF$ belong to that space:
$\mathcal G^p$ is a polynomial of degree $2sp$, and invariant
derivatives do not increase its degree.  Thus they are controlled
by finitely many seminorms of $f$ at exponent $B+2sp$.
The error estimates then put the cutoff functions and their word
images in the same pairing domain before any use of continuity.

Use \eqref{pbw-level-zero-identity},
\eqref{pbw-pairing-bound}, and the exact identity
\[
 \beta(\mathcal C_wF_\epsilon,F_\epsilon)
       -\beta(\mathcal C_wF,F)
 =\beta(\mathcal C_w(F_\epsilon-F),F_\epsilon)
       +\beta(\mathcal C_wF,F_\epsilon-F).
\]
Both terms tend to zero.  This is convergence of the original
quadratic integral, not merely of its integrand at individual
points.  Dividing by the original fixed positive word normalizer
and taking the real part proves the normalized assertion.
Finally, when $N=3s$,
\[
 m=12s^2-6s+5<12s^2=2p\qquad(s\ge1).
\]
\end{proof}

\paragraph{Combination with the normalized spectral estimate.}
If the same-form spectral calculation supplies
$E_{0,w}(F)\le jN^2(F)$ for all the allowed words, then the
finite maximum in \eqref{adaptive-all-head-bound} and
Theorem~\ref{pbw-level-zero-cutoff} give precisely
\eqref{adaptive-uncut-endpoint}, with $p=6s^2$.
For a polynomial compact-type shift $A$ of degree at most $d_A$,
apply the theorem to $f=F_0A$; its required reserve is
$A_0-d_A-2sp\ge B$.
All constants in the PBW expansion and
cutoff error are used only to establish this exact limit.  They
do not become factors in the resulting pointwise inequality.
The uniform constants still come solely from the grid scale,
the scalar all-head recursion, the normalized spectral estimate,
and the subsequent determinant moment and compact-type
comparisons.
